% Pass5396--5401: all-q Levi apartment cycle tight-frame theorem.
\subsection{All-\(q\) Levi apartments form the Hodge cycle tight frame}

Continue with a generalized quadrangle \(\mathcal S\) of order \((q,q)\), its Levi graph
\(\Gamma\), and its flag line graph \(X=L(\Gamma)\).  An apartment of the rank-two
building is a simple \(8\)-cycle in \(\Gamma\).  Let
\[
C\in\{-1,0,1\}^{N\times M},
\qquad
N=(q+1)^2(q^2+1),
\]
be the signed flag-edge/apartment incidence matrix: every Levi edge is oriented from its
point endpoint to its line endpoint, and every apartment column is given either cyclic
orientation.  Reversing an apartment reverses its entire column and therefore leaves
\(CC^{\mathsf T}\) unchanged.

The apartment counts follow directly from the generalized-quadrangle axiom.  A
nonbacktracking Levi path of five edges has endpoints consisting of a point and a
nonincident line.  The unique projection of that point to that line supplies the unique
three-edge return path, hence
\[
\boxed{\text{every nonbacktracking five-edge Levi path lies in a unique apartment}.}
\]
For a shorter fixed nonbacktracking path, the terminal vertex has \(q\) legal continuations
after excluding the edge just traversed.  Each apartment containing the shorter path chooses
exactly one of them.  Iterating therefore gives
\[
\boxed{
\#\{\text{apartments containing flag edges }e,f\}=q^{4-d_X(e,f)}.
}
\]
In particular every flag edge lies in \(q^4\) apartments.  Since every apartment contains
eight flag edges,
\[
\boxed{
M=\frac{Nq^4}{8}
=\frac{(q+1)^2(q^2+1)q^4}{8}.
}
\]

The fixed point-to-line orientation alternates signs around an apartment.  Thus two flag
edges at line-graph distance \(d\) contribute relative sign \((-1)^d\), and the full overlap
kernel is
\[
\boxed{
CC^{\mathsf T}
=q^4A_0-q^3A_1+q^2A_2-qA_3+A_4.
}
\]
Passes 5388--5395 identified the right-hand side as \(N E_{-2}\).  Therefore
\[
\boxed{
CC^{\mathsf T}=N E_{-2}=N P_{\rm cyc}.
}
\]
Consequently the complete apartment family spans the entire Levi cycle space,
\[
\boxed{\operatorname{rank}C=q^4,}
\]
and its nonzero frame-operator eigenvalue is \(N\).

Each apartment column has squared norm \(8\).  After dividing every column by \(\sqrt8\),
the \(M=Nq^4/8\) apartment vectors form a unit-norm tight frame for the
\(q^4\)-dimensional cycle space, with frame bound
\[
\boxed{N/8.}
\]
This is an exact all-\(q\) extension of the W33 cycle-frame theorem, not an analogy.

At \(q=3\),
\[
N=160,\qquad q^4=81,\qquad M=\frac{160\cdot81}{8}=1620,
\]
and
\[
\boxed{
CC^{\mathsf T}
=81A_0-27A_1+9A_2-3A_3+A_4
=160E_{-2}.
}
\]
Thus the former BT545--BT549 identities
\[
160\cdot81=1620\cdot8=12960,
\qquad
\operatorname{spec}(CC^{\mathsf T})=160^{81}\oplus0^{79},
\]
are the \(q=3\) specialization of the general apartment tight-frame law.

\paragraph{Prior-art and evidence boundary.}
The flag association scheme and its symmetric four-class fusion are existing literature;
targeted literature search also finds other tight-frame constructions built from generalized
quadrangles.  No priority claim is made here for generalized-quadrangle tight frames in the
broad sense.  The theorem asserted here is the explicit full-apartment cycle-incidence identity
above and its Hodge specialization.  It does not assert tightness for arbitrary apartment
subsets and does not by itself imply a code distance or physical fault-tolerance threshold.
